\ChapterImageStar[cap:justificacion]{Justificación}{./images/fondo.png}\label{cap:justificacion}
\mbox{}\\

\noindent
La ciberseguridad según~\cite{lallie2025universitycyber} se ha convertido en un tema que ha cobrado creciente relevancia 
en un mundo donde la tecnología abarca diversos ámbitos laborales a nivel mundial, debido al aumento de incidentes que comprometen 
la integridad de la información de organizaciones tanto públicas como privadas. Este panorama también contiene a las instituciones 
de educación superior, las cuales requieren de infraestructuras tecnológicas con la capacidad de soportar y proteger sus activos de 
información~\cite{educause2023cyberhighered}. En contexto de la Universidad del Quindío, este trabajo de grado enfatiza en el grupo 
de investigación GRID, debido a que éste dispone de recursos tecnológicos estratégicos desarrollados con el fin de contener activos 
importantes para el grupo, como información, investigaciones y ofertas académicas dispuestas para los estudiantes, sin embargo, no 
se evidencian controles adaptados a estándares internacionales de la familia de las ISO/IEC 27000, igualmente, el UTM SonicWall NSA 
6600 activo de hardware para la seguridad perimetral, evidencia actualmente limitaciones en la prestación efectiva de sus servicios 
de seguridad para asumir dicho rol dentro de la infraestructura de computación en la nube del grupo GRID.\@ Por lo anterior, su 
arquitectura de seguridad requiere un fortalecimiento que permita enfrentar de manera sistemática las amenazas 
emergentes~\cite{edutech2025studentdata}. Bajo este escenario, la propuesta se centra en especificar una arquitectura de seguridad 
para la infraestructura de la computación en la nube del grupo de investigación GRID, lo que busca contribuir con la administración 
institucional de la ciberseguridad y la implementación de un determinado conjunto de controles establecidos en la familia ISO/IEC 27000.
